% Appendix
\appendix

\chapter{Appendix: Resources and Licensing}
\label{app:resources}

\section{Deployment and Access}
\label{app:deployment}

The interactive Streamlit application is deployed and accessible online:

\begin{itemize}
  \item \textbf{Live Deployment URL:} http://134.199.172.48/
  \item \textbf{Application Type:} Streamlit web app
  \item \textbf{Infrastructure:} DigitalOcean \$4/month VPS (Ubuntu 24.04.2 LTS)
  \item \textbf{Port:} 8501 (reverse-proxied via caddy)
\end{itemize}

Users can:
\begin{enumerate}
  \item Select and analyze existing Pokémon from the dataset.
  \item Create and predict custom Pokémon specifications.
  \item View model metrics, feature importance rankings, and engineered feature snapshots.
  \item Experiment with different stat combinations to understand model decision boundaries.
\end{enumerate}

\section{Source Code Repository}
\label{app:repository}

All source code, data, and model artifacts are available on GitHub:

\begin{itemize}
  \item \textbf{Repository URL:} \url{https://github.com/shravanasati/unown}
  \item \textbf{Repository Name:} \texttt{unown}
  \item \textbf{Owner:} \texttt{shravanasati}
  \item \textbf{Primary Branch:} \texttt{main}
\end{itemize}

\subsection{Installation and usage}

To run the project locally:

\begin{verbatim}
# Clone the repository
git clone https://github.com/shravanasati/unown.git
cd unown

# Install dependencies
uv sync

# Run the Streamlit app
uv run streamlit run streamlit_app.py

# Run model training (if retraining desired)
jupyter notebook model_training.ipynb
\end{verbatim}

\section{License}
\label{app:license}

This project is licensed under the \textbf{Revised BSD License} (BSD-3-Clause). 

\subsection{Third-party dependencies}

The project depends on the following libraries, each governed by its own license.

Users are responsible for complying with all applicable licenses when using or distributing this software and its dependencies.

\section{Data Attribution}
\label{app:data_attribution}

The Pokémon dataset is sourced from the official Pokémon franchise game specifications, accessed via the \href{https://pokeapi.co}{\textbf{PokeAPI}}. The data includes:
\begin{itemize}
  \item Base statistics (HP, Attack, Defense, Sp. Atk, Sp. Def, Speed)
  \item Pokédex information (name, generation, type, abilities, height, weight, base experience)
  \item Legendary status (manually curated based on official Pokémon classifications)
\end{itemize}

Pokémon is a trademark of The Pokémon Company, Nintendo, and Game Freak. This project is for educational and research purposes only.

\section{Citation}
\label{app:citation}

If you use this project in academic work, research, or publications, please cite it as follows:

\begin{verbatim}
@software{unown2025,
  author = {Shravan Asati},
  title = {Unown: Pokemon Legendary Status Predictor},
  url = {https://github.com/shravanasati/unown},
  year = {2025},
  license = {BSD-3-Clause}
}
\end{verbatim}

Or in BibTeX format:

\begin{verbatim}
@inproceedings{unown2025,
  title={Domain-Informed Feature Engineering for Imbalanced Classification: 
         A Pokemon Legendary Detection Case Study},
  author={Asati, Shravan},
  year={2025},
  note={Available at https://github.com/shravanasati/unown}
}
\end{verbatim}

\section{Contact and Support}
\label{app:contact}

For questions, bug reports, or contributions:

\begin{itemize}
  \item \textbf{GitHub Issues:} \url{https://github.com/shravanasati/unown/issues}
  \item \textbf{GitHub Discussions:} \url{https://github.com/shravanasati/unown/discussions}
  \item \textbf{Pull Requests:} Contributions are welcome; see repository for guidelines.
\end{itemize}
